\paragraph{Exhaustive Boolean-domain certificates.}
A separate, predeclared cohort trains width-$64$ tanh networks on $400$ of the $512$ vertices of $\{-1,1\}^9$, with the remaining $112$ vertices held out.  For the protocol-selected depth-$4$ majority network \texttt{500}, exact arithmetic and outward intervals certify the clipped-score condition on \emph{every} Boolean input at $c=.03$.  All $400$ training margins are independently interval-verified above $1/2$, and the network correctly predicts all $512$ labels.  The global raw-change bound is $.162903>1/16$, whereas every clean absolute output exceeds $4.72037$.  Consequently the clipped score is exactly constant under every permitted disturbance at every vertex.  A separate auditor replayed the saved matrix factors, all $512$ interval trajectories, and all $400$ training inequalities.

A subsequent, separately recorded check gives the same exhaustive conclusion for the depth-$16$ majority network \texttt{501} at the largest primary radius $c=1/16$.  Its raw-change bound is $3.474718$, while every clean absolute output exceeds $5.12034$; all $512$ labels and every training margin are correct.  These finite-domain certificates use the elementary implication
\[
 |z(x)|\ge m,\qquad |z^\eta(x)-z(x)|\le B\le m-\tfrac12
 \quad\Longrightarrow\quad
 \operatorname{clip}_{1/2}(z^\eta(x))=\operatorname{clip}_{1/2}(z(x)).
\]
They therefore certify the actual clipped-score condition even when the stronger raw-output-change condition is not certified.  Their scope is the entire finite Boolean domain, not all real inputs.

\paragraph{A depth-$64$ example requiring reachable-state bounds.}
The first depth-$64$ majority network, \texttt{504}, also satisfies the condition on \emph{all} $512$ Boolean vertices at the primary $c=.001$.  Every one of the $400$ training margins is interval-verified above $7.61170$, and every clean absolute output on the full cube exceeds $5.64786$.  Thus this network also satisfies the original all-input clean-margin restriction, not only the newer clipped-score restriction.  Exact interval label comparisons give $511/512$ correct predictions on the full cube, including $111/112$ held-out predictions; the remaining vertex is a confidently wrong prediction, not an omitted or ambiguous case.  Its global raw-change bound, $13.414533$, cannot establish this conclusion, even with output clipping.  Instead, $52$ finite groups of vertices share rigorously bounded reachable-state tubes.  Each group's disturbed outputs remain beyond the corresponding clipping threshold, with the smallest verified surplus exceeding $1.40770$.  Thus the clipped score is exactly invariant for every allowed disturbance at every vertex.  This is an exhaustive finite-domain certificate for an ordinary SGD-trained depth-$64$ network, rather than a sampled-input audit or a claim about every real input.

For completeness, sharing a tube across finitely many inputs does not alter the underlying proof.  At each layer and coordinate, replace the clean preactivation's distance from zero by the minimum of that distance over all vertices in the group.  A common prefix-energy bound then encloses every group's noisy trajectory, so the same derivative and matrix-norm bounds apply to each vertex.  The grouping only reduces computation.  Every input is still checked, and any unresolved input would prevent a full-domain conclusion.  The certificate uses exact integer matrix arithmetic, outward interval $\tanh$, saved norm-factor proposals, and explicit source hashes; it is not a Lean proof.  A separate auditor regenerated all $512$ clean interval trajectories, checked all $52$ groups and $3{,}216$ saved matrix bounds, and replayed every training inequality.  The three worst groups were additionally checked with the original pure-integer matrix-norm verifier.

\paragraph{A depth-$64$ nonlinear target.}
The corresponding recursive-majority network \texttt{505}, trained on majority of three disjoint three-bit majorities, also satisfies the clipped-score condition on the entire $512$-vertex cube at $c=.001$.  Every one of its $400$ training margins is interval-verified above $7.47388$, and every clean absolute output exceeds $2.40649$, so the original all-input clean-margin condition holds as well.  Its full-domain accuracy is $504/512$; the held-out accuracy is $104/112$.  All eight errors remain in the verified domain.

For this network, the shared reachable-tube proof establishes $511$ of the $512$ cases.  The remaining vertex, index $325$, is explicitly retained as unresolved in that component certificate.  A separate finite-radius proof propagates the exact clean-response covariance and bounds the nonlinear Taylor remainder in state norm.  It encloses the noisy output at this same vertex in $[-2.586691,-2.226290]$, proving that its clipped score is also invariant.  The composite certificate records the union of these two proofs, with model, radius, source, and component hashes; the input domain never changes.  The exceptional-vertex proof was independently replayed, including all $126$ exact matrix-norm bounds.  A further independent replay regenerated all $512$ clean trajectories, all $122$ groups and $7{,}585$ local matrix bounds, and every training margin, then checked the composite coverage.  The state-covariance inequality used for this refinement is given below.
